\begin{frame}{MDP의 다섯 요소}
  MDP는 다섯 가지로 정의된다: \kb{$(\mathcal{S}, \mathcal{A}, P, R,
  \gamma)$}
  \small
  \begin{tabular}{@{}ll@{}}
    \toprule
    요소를 한 줄로 & 의미 \\
    \midrule
    $\mathcal{S}$ & 가능한 \kb{상태}들의 집합 \\
    $\mathcal{A}$ & 가능한 \kb{행동}들의 집합 \\
    $P(s'|s,a)$ & $s$에서 $a$를 했을 때 $s'$로 \kb{전이}될 확률 \\
    $R(s,a)$ & $s$에서 $a$를 했을 때 받는 \kb{즉시 보상} \\
    $\gamma \in [0,1)$ & \kb{할인율} (discount factor) \\
    \bottomrule
  \end{tabular}
  \vskip1em
  \begin{block}{Week 2의 밴딧과 비교}
    밴딧에는 $\mathcal{S}$도 $P$도 \textbf{없었다}(팔을 당겨도
    ``다음 상태''가 없었다). MDP는 지금 한 행동이 다음 상태
    $P(s'|s,a)$ 를 통해 \textbf{미래 전체}에 영향을 준다는 사실을
    명시적으로 담는다.
  \end{block}
\end{frame}
